% Slide build of Figure 28.1 -- the reachability ladder, one rung at a time.
%
% Printed whole, the ladder is a reference card: the reader scans it and moves
% on. Built, it is a method being demonstrated, which is what a lecture is for.
%
% But this figure must NOT be revealed by accumulation. Its shape -- five rungs,
% each eliminating everything below -- is itself the lesson, and a build that
% adds one rung at a time spends the first three steps showing a ladder that
% appears to have three rungs, with a void where the rest of the method should
% be. So every rung is drawn faint from the start and lights up as it is
% reached: the class can see how far there is to go, which is exactly the
% reassurance the method is meant to give a person staring at a broken laptop.
%
% \stepdim{n}{..} is the complement of \stepvis{n}{..} -- it renders only
% while n is still ahead. A dim/solid pair draws the same element twice, once
% in each state, and never both.
\begin{tikzpicture}[font=\scriptsize]
  \useasboundingbox (-2.2,-7.95) rectangle (9.6,0.9);

  \tikzset{
    rung/.style={rounded corners=3pt, draw=#1, fill=#1!8, text=#1!45!black,
                 line width=0.9pt, text width=3.5cm, align=left,
                 minimum height=1.0cm, font=\scriptsize, inner sep=5pt},
    ghost/.style={rounded corners=3pt, draw=neutral!30, fill=neutral!3,
                  text=neutral!35, line width=0.9pt, text width=3.5cm,
                  align=left, minimum height=1.0cm, font=\scriptsize,
                  inner sep=5pt},
    res/.style={font=\tiny, text=#1!45!black, align=left, anchor=west},
    fl/.style={-{Stealth[length=2.2mm]}, line width=0.9pt, neutral!70},
    stepcap/.style={font=\bfseries, anchor=north west, text width=11.3cm,
                    align=left}}

  % Rung 1 is live from the first beat; 2-5 start as ghosts.
  \node[rung=dom1] at (0,0)
      {\textbf{1.} Do I have an address?\\ \cmd{ipconfig /all}};

  \stepdim{2}{\node[ghost] at (0,-1.5) {\textbf{2.} Ping my own gateway};}
  \stepvis{2}{\node[rung=dom1] at (0,-1.5) {\textbf{2.} Ping my own gateway};}

  \stepdim{3}{\node[ghost] at (0,-3.0)
      {\textbf{3.} Ping something remote\\ \emph{by IP address}};}
  \stepvis{3}{\node[rung=dom1] at (0,-3.0)
      {\textbf{3.} Ping something remote\\ \emph{by IP address}};}

  \stepdim{4}{\node[ghost] at (0,-4.5)
      {\textbf{4.} Ping the same thing\\ \emph{by name}};}
  \stepvis{4}{\node[rung=dom1] at (0,-4.5)
      {\textbf{4.} Ping the same thing\\ \emph{by name}};}

  \stepdim{5}{\node[ghost] at (0,-6.0) {\textbf{5.} It is not the network};}
  \stepvis{5}{\node[rung=dom2] at (0,-6.0)
      {\textbf{5.} It is not the network};}

  % Coordinates, not node names -- the nodes are drawn conditionally, so a
  % \draw naming one would be an undefined-node error at the wrong level.
  \draw[fl, neutral!25] (0,-0.55) -- (0,-0.95);
  \draw[fl, neutral!25] (0,-2.05) -- (0,-2.45);
  \draw[fl, neutral!25] (0,-3.55) -- (0,-3.95);
  \draw[fl, neutral!25] (0,-5.05) -- (0,-5.45);
  \stepvis{2}{\draw[fl] (0,-0.55) -- (0,-0.95);}
  \stepvis{3}{\draw[fl] (0,-2.05) -- (0,-2.45);}
  \stepvis{4}{\draw[fl] (0,-3.55) -- (0,-3.95);}
  \stepvis{5}{\draw[fl] (0,-5.05) -- (0,-5.45);}

  % The verdicts are the payload: each arrives with the rung it belongs to,
  % and unlike the rungs there is no ghost -- an unread verdict would be a
  % spoiler, and the question "so what does it mean if this one fails?" is
  % the one worth asking the room before advancing.
  \stepvis{1}{\node[res=accent] at (2.4,-0.75)
      {fails $\rightarrow$ \textbf{DHCP or layer 2}\\Chapters 20, 14};}
  \stepvis{2}{\node[res=accent] at (2.4,-2.25)
      {fails $\rightarrow$ \textbf{local segment}\\VLAN, cable, ARP,
       wireless assoc.};}
  \stepvis{3}{\node[res=accent] at (2.4,-3.75)
      {fails $\rightarrow$ \textbf{routing, NAT or ACL}\\Chapters 15, 22, 24};}
  \stepvis{4}{\node[res=accent] at (2.4,-5.25)
      {fails $\rightarrow$ \textbf{DNS only}\\Chapter 21};}

  \steponly{1}{\node[stepcap, text=dom1!45!black] at (-2.0,-6.7)
      {Start at the host, not at the core. No address, and nothing else in
       this list can possibly work.};}
  \steponly{2}{\node[stepcap, text=dom1!45!black] at (-2.0,-6.7)
      {The gateway is the last hop still on our own segment. Passing here
       clears the cable, the VLAN and the association.};}
  \steponly{3}{\node[stepcap, text=dom1!45!black] at (-2.0,-6.7)
      {By address, on purpose --- it takes name resolution out of the
       question entirely.};}
  \steponly{4}{\node[stepcap, text=dom1!45!black] at (-2.0,-6.7)
      {Only now bring names back. If 3 passed and 4 fails, exactly one thing
       is left, and it is not the network.};}
  \steponly{5}{\node[stepcap, text=dom2!45!black] at (-2.0,-6.7)
      {Every rung that passes eliminates everything below it. Four commands,
       and the fault is in one of five places.};}
\end{tikzpicture}
